Throughout this article, rings have an identity and modules are unital left modules. The ring need not be commutative. The ring-theoretic ascending chain condition is introduced in \wikilink{s29--b001}{the Noetherian ring definition}; here it is applied to all submodules of a module. Quotients and exact sequences are as in \wikilink{a2s08--b009}{the module isomorphism theorems} and \wikilink{a2s09--b007}{short exact sequences}.

\begin{definition}\label{def-finitely-generated-module}
An $R$-module $M$ is \emph{finitely generated} if there are $m_1,\ldots,m_n\in M$ such that $M=Rm_1+\cdots+Rm_n$. Equivalently, there is a surjective $R$-linear map $R^n\to M$ for a finite nonnegative integer $n$.
\end{definition}

\begin{definition}\label{def-noetherian-module}
An $R$-module $M$ is \emph{Noetherian} if every ascending sequence of submodules
\[
N_1\subseteq N_2\subseteq N_3\subseteq\cdots\subseteq M
\]
stabilizes: there is an $n$ such that $N_j=N_n$ for all $j\ge n$. This is the \emph{ascending chain condition on submodules}.
\end{definition}

\begin{proposition}\label{prop-noetherian-equivalences}
For an $R$-module $M$, the following conditions are equivalent.
\begin{enumerate}
\item $M$ is Noetherian.
\item Every submodule of $M$ is finitely generated.
\item Every nonempty family of submodules of $M$ has a maximal member under inclusion.
\end{enumerate}
\end{proposition}
\begin{proof}\label{proof-noetherian-equivalences}
Suppose that $M$ satisfies the ascending chain condition. If a submodule $N$ were not finitely generated, choose $x_1\in N$, then $x_2\in N\setminus Rx_1$, and inductively $x_{j+1}\notin Rx_1+\cdots+Rx_j$. This gives a strictly ascending sequence, a contradiction. Thus every submodule is finitely generated.

Conversely, let $N_1\subseteq N_2\subseteq\cdots$ be an ascending sequence and set $N=\bigcup_jN_j$. This union is a submodule. If $N$ has finitely many generators, all of them lie in some $N_n$, so $N=N_n$ and the sequence stabilizes.

If the ascending chain condition holds and a nonempty family has no maximal member, start with a member and repeatedly choose a strictly larger member. The resulting sequence contradicts the condition. Finally, if every nonempty family has a maximal member, apply this to the members of any ascending sequence. Once a maximal member occurs, all later members equal it. These arguments use the usual choice convention for making successive selections.
\end{proof}

\begin{theorem}\label{thm-noetherian-exact-sequences}
Let
\[
0\longrightarrow N\xrightarrow{i}M\xrightarrow{q}Q\longrightarrow0
\]
be an exact sequence of left $R$-modules. Then $M$ is Noetherian if and only if both $N$ and $Q$ are Noetherian.
\end{theorem}
\begin{proof}\label{proof-noetherian-exact-sequences}
Identify $N$ with its image in $M$. If $M$ is Noetherian, every ascending chain in $N$ is one in $M$. An ascending chain in $Q$ lifts by inverse image under $q$ to an ascending chain in $M$, so it too stabilizes. Thus $N$ and $Q$ are Noetherian.

Conversely, suppose that $N$ and $Q$ are Noetherian, and let $L\subseteq M$ be a submodule. The submodule $L\cap N$ is finitely generated, say by $u_1,\ldots,u_r$. The submodule $q(L)$ of $Q$ is generated by finitely many elements $q(v_1),\ldots,q(v_s)$ with $v_j\in L$. For any $x\in L$, write $q(x)=\sum_ja_jq(v_j)$. Then $x-\sum_ja_jv_j\in L\cap N$, so $L$ is generated by the $u_i$ and $v_j$. Every submodule of $M$ is therefore finitely generated, proving that $M$ is Noetherian.
\end{proof}

\begin{corollary}\label{cor-finite-sums-noetherian}
A finite direct sum of Noetherian modules is Noetherian. Submodules and quotients of a Noetherian module are Noetherian.
\end{corollary}
\begin{proof}\label{proof-finite-sums-noetherian}
The submodule and quotient assertions are contained in the theorem. For a direct sum of two modules, apply it to the split exact sequence $0\to M_1\to M_1\oplus M_2\to M_2\to0$. Induction gives the assertion for any finite number of summands, including the zero module.
\end{proof}

\begin{theorem}\label{thm-finite-modules-noetherian-ring}
If $R$ is left Noetherian, every finitely generated left $R$-module is Noetherian. In particular, every submodule of a finitely generated left $R$-module is finitely generated.
\end{theorem}
\begin{proof}\label{proof-finite-modules-noetherian-ring}
Submodules of the left regular module ${}_RR$ are exactly left ideals. Thus the definition of a left Noetherian ring says that ${}_RR$ is a Noetherian module. A finite direct sum $R^n$ is Noetherian by the corollary. Every finitely generated module is a quotient of such an $R^n$, hence is Noetherian. The finite generation of its submodules follows from the equivalence proved above.
\end{proof}

This explains the Noetherian hypothesis behind \wikilink{a2s10--b008}{the category of finitely generated modules}: kernels are submodules and remain finitely generated, while cokernels are finitely generated quotients. The source phrase ``Noetherian category'' there is understood as a Noetherian ring. The \wikilink{a2s15--b005}{corresponding PID result} is a special case, not a restriction on the general theorem.

\begin{example}\label{ex-infinite-sum-not-noetherian}
The free abelian group $\bigoplus_{j\ge1}\mathbb Z e_j$ is not Noetherian. The subgroups $\mathbb Ze_1+\cdots+\mathbb Ze_n$ form a strictly ascending sequence. Thus an infinite direct sum of Noetherian modules need not be Noetherian, and finite generation of a module over an arbitrary ring is not by itself the Noetherian condition.
\end{example}

\begin{definition}\label{def-hopfian-module}
An $R$-module $M$ is \emph{Hopfian} if every surjective $R$-linear endomorphism of $M$ is an automorphism.
\end{definition}

\begin{theorem}\label{thm-noetherian-hopfian}
Every Noetherian module is Hopfian. Consequently, every finitely generated module over a left Noetherian ring is Hopfian.
\end{theorem}
\begin{proof}\label{proof-noetherian-hopfian}
Let $f:M\to M$ be surjective, with $M$ Noetherian. The ascending sequence $\ker f\subseteq\ker f^2\subseteq\cdots$ stabilizes, so for some $n\ge1$ we have $\ker f^n=\ker f^{n+1}$. If $x\in\ker f$, surjectivity of $f^n$ gives $y\in M$ with $f^n(y)=x$. Then $f^{n+1}(y)=f(x)=0$, so $y\in\ker f^{n+1}=\ker f^n$, and $x=0$. Thus $f$ is injective as well as surjective, and its inverse is $R$-linear. The final assertion follows from the preceding theorem.
\end{proof}

\begin{remark}\label{remark-hopfian-not-noetherian}
The converse is false. As a $\mathbb Z$-module, $\mathbb Q$ is not finitely generated and hence is not Noetherian. Every additive endomorphism $f$ of $\mathbb Q$ is multiplication by $f(1)$: for integers $a$ and $b\ne0$, the identity $b f(a/b)=a f(1)$ determines $f(a/b)$. A surjective such endomorphism has $f(1)\ne0$ and is invertible. Thus $\mathbb Q$ is Hopfian.
\end{remark}

\begin{definition}\label{def-directly-finite-ring}
A unital ring $R$ is \emph{directly finite} if $ab=1$ implies $ba=1$ for all $a,b\in R$.
\end{definition}

\begin{corollary}\label{cor-left-noetherian-directly-finite}
Every left Noetherian ring is directly finite.
\end{corollary}
\begin{proof}\label{proof-left-noetherian-directly-finite}
Suppose $ab=1$. For $c\in R$, let $\rho_c:R\to R$ be right multiplication, $\rho_c(x)=xc$. This is an endomorphism of the left regular module, even when $R$ is not commutative. The identity $ab=1$ gives $\rho_b\rho_a=\mathrm{id}_R$, so $\rho_b$ is surjective. The left regular module is Noetherian, and therefore $\rho_b$ is injective by the Hopfian theorem. For any $x\in R$,
\[
\rho_b\bigl(\rho_a\rho_b(x)\bigr)=\rho_b(x).
\]
Injectivity yields $\rho_a\rho_b(x)=x$. Taking $x=1$ gives $ba=1$.
\end{proof}
